\documentclass[final,3p,times]{elsarticle}

\usepackage[UTF8]{ctex}
\usepackage{graphicx}
\usepackage{booktabs}
\usepackage{amsmath,amssymb}
\usepackage{siunitx}
\usepackage{enumitem}
\usepackage{microtype}
\usepackage[hidelinks]{hyperref}
\usepackage{cleveref}
\usepackage{titlesec}

\setlist[itemize]{leftmargin=2em,topsep=0.35em,itemsep=0.25em,parsep=0em}
\setlist[enumerate]{leftmargin=2em,topsep=0.35em,itemsep=0.25em,parsep=0em}

\titleformat{\section}{\normalfont\bfseries\large}{\thesection}{0.8em}{}
\titleformat{\subsection}{\normalfont\bfseries\normalsize}{\thesubsection}{0.8em}{}
\titleformat{\subsubsection}{\normalfont\bfseries\normalsize}{\thesubsubsection}{0.8em}{}
\titlespacing*{\section}{0pt}{1.1em}{0.6em}
\titlespacing*{\subsection}{0pt}{0.9em}{0.4em}
\titlespacing*{\subsubsection}{0pt}{0.8em}{0.3em}

\crefname{table}{表}{表}
\crefname{figure}{图}{图}
\Crefname{table}{表}{表}
\Crefname{figure}{图}{图}

\begin{document}

\begin{frontmatter}

\title{Kestrel：基于视觉-语言大模型的复杂曲面OCR驱动PIVAS药品配液精准校验系统}

\author[inst1]{作者信息待补充}
\address[inst1]{Beijing Information Science and Technology University}

\begin{abstract}
静脉用药调配中心（PIVAS）在配液环节需要对药品名称、规格、批号、有效期、剂量与溶媒等关键信息进行终末核对。现行人工校验效率受限且易疲劳；在复杂曲面药瓶（如圆柱体）场景下，标签遮挡、反光与起皱等因素进一步导致传统OCR识别准确率下降难以落地。
本文设计一种基于小参数视觉-语言大模型（VLM）的药品配液精准校验系统：系统将前端采集视频经降采样与预处理后输入VLM，完成多模态OCR识别与字段抽取；VLM输出包含药品名称、规格、批号、有效期、剂量与溶媒等字段的结构化信息，并通过知识库约束与业务规则实现规范化对齐与一致性校验。
在真实或准真实PIVAS数据集与复杂药瓶难例上，系统在遮挡与标签退化等条件下仍能保持稳定的信息提取与校验流程；相较传统OCR流水线，基于VLM的识别方式在效率与准确度上表现出显著优势。
本文系统为PIVAS配液校验提供了可部署、可追溯且稳健的技术方案，其关键创新点在于以VL大模型替代传统OCR作为核心识别引擎，并结合知识库约束降低不确定性与错误传播风险。
\end{abstract}

\begin{keyword}
\normalfont PIVAS \sep 视觉-语言大模型 \sep 复杂曲面OCR \sep 多模态OCR \sep 知识库约束
\end{keyword}

\end{frontmatter}

\section{引言}
静脉用药调配中心（PIVAS）承担大量静脉用药的集中配制任务，流程通常包括医嘱审核、药品拣选、配液操作、成品贴签与终末校验等关键环节。其中终末校验直接关系到患者用药安全，需重点识别并防范如“用错药、用错规格、用错批号/效期、剂量/溶媒不符、成品标签不一致”等差错类型。

目前，PIVAS校验实践往往依赖人工目视复核、条码扫描以及部分规则检查。人工复核在高强度重复作业下容易出现注意力下降与疲劳累积，效率和稳定性难以保证；条码方案在条码缺失、破损或被遮挡时会失效；而传统规则匹配高度依赖上游识别的准确性，一旦关键字段识别受噪声影响，容易导致误报或漏报。更重要的是，药瓶多为圆柱体且材质易反光，标签在搬运与存放过程中可能出现磨损、起皱、污染或贴附不平整，再叠加拍摄角度变化与运动模糊，导致传统OCR在复杂药瓶场景下性能显著退化。

本文关注的核心空白在于：在真实PIVAS操作约束（低延迟、可追溯、可部署、可与药品主数据/医嘱系统对接）下，如何解决“药瓶标签OCR稳定性不足”这一瓶颈，并将识别结果可靠地用于校验决策。传统OCR往往以字符级准确为目标，但在复杂药瓶条件下存在三个典型问题：（1）字符不可见或严重失真导致漏识；（2）识别结果碎片化且难以对齐到业务字段（例如药名、规格、批号、有效期）；（3）相似包装与缩写表达导致字符串匹配误判。与之相比，视觉-语言大模型具备跨模态语义理解能力，能够在不完全文字可见的情况下，结合上下文线索完成字段级解析，并输出结构化结果与不确定性标记，从而更适合作为PIVAS校验的“识别引擎”。

为此，我们提出一种基于小参数视觉-语言大模型的药品配液精准校验系统，以多模态OCR与字段理解为核心，结合知识库约束与规则校验实现保守可控的校验输出。本文主要贡献如下：
\begin{itemize}
  \item 提出一种以视觉-语言大模型为核心识别引擎的PIVAS药品配液精准校验系统，输出结构化字段、差错分型与可追溯证据。
  \item 面向复杂药瓶OCR难题，提出“多帧证据聚合+VLM字段级抽取”的识别策略，提升遮挡、反光与磨损条件下的鲁棒性。
  \item 引入知识库约束与规则校验，将识别结果规范化对齐到药品主数据与医嘱上下文，并在证据不足时触发保守复核，降低错误传播风险。
\end{itemize}

\section{方法}
\subsection{系统总体流程}
给定调配台面图像或短视频输入，系统输出三类结果：（1）校验结论（通过/拦截/人工复核）；（2）差错类型（例如药品不符、规格不符、批号或效期异常、剂量/溶媒不符、标签不一致等）；（3）可追溯证据（定位区域、可见文本片段、规范化字段、触发的规则与置信度）。整体流程由四个阶段构成：定位、OCR、多模态推理与规则校验。

\subsection{采集层与中间层（简述）}
系统采用“前端采集--中间层通信--后端推理与校验”的分布式形态。前端通过多视角/多帧采集为识别提供更充分的视觉证据；中间层负责将图像数据与控制信令可靠传输到后端，并维护任务状态与日志追溯所需的基础信息。本文重点不在采集与通信实现细节，而在后端识别与校验算法设计。

\subsection{定位与多帧视觉证据聚合}
在每一帧上执行药瓶与标签区域定位，并对候选区域进行质量评估（清晰度、可读性、反光与遮挡程度等）。系统从序列中选择若干高质量帧进行聚合：一方面对可见文本片段进行去重与一致化，另一方面为后续VLM提取提供更充分的视觉证据，从而提升复杂药瓶场景下的稳定性。

\subsection{复杂药瓶OCR难点}
PIVAS药瓶标签识别的主要困难并非单一“识别字符”，而是“在复杂成像条件下稳定获取可用于业务校验的字段”。具体而言：（1）圆柱体曲面导致文字形变与局部不可见，且手部/相邻药瓶易造成遮挡；（2）瓶体反光与标签覆膜产生高亮区域，导致字符断裂或对比度不足；（3）标签磨损、污染或起皱引入随机噪声，使字符边界不清；（4）不同厂家与批次版式差异大，模板依赖方法维护成本高；（5）药名与规格存在缩写、别名与单位表达差异，单纯字符串匹配难以鲁棒消歧。上述因素共同导致传统OCR流水线常出现“碎片化输出、字段错配、低置信度难以决策”等问题。

\subsection{基于小参数视觉-语言大模型的多模态OCR与结构化抽取}
本文将小参数视觉-语言大模型作为多模态识别与理解核心：在给定标签区域与多帧证据的条件下，模型直接完成“读字（OCR）+字段理解（语义解析与规范化）”的一体化输出。与传统OCR先识别字符再做规则拼接不同，VL模型能够在更高的语义层面进行聚合：利用上下文将碎片化文字对齐到业务字段，并对不完整证据输出不确定性标记。通过提示词约束固定输出结构，字段示例如下：
\begin{enumerate}
  \item 药品名称（通用名/商品名）；
  \item 规格信息（含剂量、体积、浓度等表达）；
  \item 批号；
  \item 有效期；
  \item 生产厂家（可选）；
  \item 不确定性标记与证据引用。
\end{enumerate}
为保持系统保守与可控，要求模型对每个字段给出来自OCR候选或可见文本片段的支撑证据；当证据不足、冲突或可见性过低时，将字段标记为未知并触发人工复核，从流程上避免“猜测式补全”带来的风险。

\subsection{知识对齐与规则校验}
结构化字段将与药品主数据进行对齐，以处理别名、规格表达差异、单位换算与相似药品消歧等问题。随后引入业务规则对字段与上下文进行一致性校验，例如：（1）药品与规格是否与医嘱一致；（2）配液时间点的有效期是否满足要求；（3）剂量与溶媒是否满足协议与约束；（4）高警示药品是否采用更严格的阈值与复核策略。规则引擎输出差错分型与风险等级，并与证据一起用于告警与追溯。

\section{结果}
\subsection{定量对比}
\Cref{tab:main-results}给出主要定量结果模板。将“待填”替换为最终指标，并按需要补充更多基线方法。

\begin{table}[htbp]
\centering
\caption{PIVAS校验基准上的主要定量对比（占位）。}
\label{tab:main-results}
\begin{tabular}{lcccc}
\toprule
方法 & 字段准确率$\uparrow$ & 校验正确率$\uparrow$ & 误报率$\downarrow$ & 时延(ms)$\downarrow$ \\
\midrule
仅OCR+字符串匹配 & 待填 & 待填 & 待填 & 待填 \\
检测+OCR+规则 & 待填 & 待填 & 待填 & 待填 \\
本文方法（VL大模型+规则） & 待填 & 待填 & 待填 & 待填 \\
\bottomrule
\end{tabular}
\end{table}

\subsection{消融实验}
\Cref{tab:ablation}给出消融模板，常见组件包括多帧聚合、知识对齐与规则约束推理等。

\begin{table}[htbp]
\centering
\caption{消融实验（占位）。}
\label{tab:ablation}
\begin{tabular}{lcc}
\toprule
设置 & 校验正确率$\uparrow$ & 说明 \\
\midrule
完整系统 & 待填 & --- \\
去掉多帧聚合 & 待填 & --- \\
去掉知识对齐 & 待填 & --- \\
去掉规则约束 & 待填 & --- \\
\bottomrule
\end{tabular}
\end{table}

\subsection{定性可视化}
\Cref{fig:qualitative}提供一个无需外部图片也可编译的占位图示例。后续可替换为真实案例，如遮挡、反光、磨损与相似包装等。

\begin{figure}[htbp]
\centering
\fbox{\parbox[c][4cm][c]{0.9\linewidth}{\centering 定性可视化占位图。\\后续替换为真实案例与说明文字。}}
\caption{带证据的定性示例：标签定位区域、OCR候选、规范化字段与触发规则（占位）。}
\label{fig:qualitative}
\end{figure}

\section{讨论}
从占位实验模板的设计出发，本文方法的优势预期主要体现在两方面：其一，以小参数视觉-语言大模型承担多模态OCR与字段理解任务，能够将碎片化、部分可见的文字线索提升到字段级语义解析，并输出不确定性标记，从而降低传统OCR“字符可见性不足即失败”的脆弱性；其二，知识库约束与规则校验为模型输出提供一致性边界，使系统在证据不足时更倾向于保守决策（触发复核而非猜测），以降低潜在风险。

本文仍存在若干局限：（1）系统效果与相机布置和图像质量相关，极端反光或严重标签缺失场景仍可能需要人工复核；（2）跨院区泛化可能受药品包装、供应商与标注规范差异影响；（3）工程部署需持续监控误报与漏报，尤其对于高警示药品需要更谨慎的阈值策略与回退机制。

未来工作可包括：扩展数据覆盖更多院区与包装形态，结合主动学习进行增量改进，探索隐私保护的部署与训练策略，并开展人因研究量化工作负担下降与可用性提升。

\section{结论}
本文面向PIVAS配液终末校验场景，提出一种以视觉-语言大模型为核心识别引擎的药品配液精准校验系统框架，将多模态OCR与字段语义理解一体化，并结合知识对齐与规则校验输出可追溯证据与保守可控的校验结论。该框架有望在保证安全性的同时提升校验效率，为智慧药房质控提供可落地的技术路径。

\section*{致谢}
本研究得到待填基金/项目名称支持。感谢待填合作单位/团队成员在数据采集与流程讨论方面提供的帮助。感谢硬件支持（如AMD AI MAX 395，若适用）。

\begin{thebibliography}{99}
  \bibitem{ref1} 作者. 题目. 期刊名, 年, 卷(期): 页码.
  \bibitem{ref2} 作者. 题目. 期刊名, 年, 卷(期): 页码.
  \bibitem{ref3} 作者. 题目. 会议名, 年: 页码.
  \bibitem{ref4} 作者. 题目. 期刊名, 年, 卷(期): 页码.
  \bibitem{ref5} 作者. 题目. 期刊名, 年, 卷(期): 页码.
\end{thebibliography}

\end{document}

